% 摘要文字、标题及关键词对应 9.12中午_误差分析完成.docx 的摘要页。
\begin{abstract}
药材烘干是中药材加工中的重要环节，合理掌握烘干过程中温度和水分的变化规律，有助于判断药材的干燥状态和所需烘干时间。本文针对圆柱形药材的烘干过程，通过分析药材内部的径向传热传质规律以及药材与烘房环境之间的交换关系，分别建立不同条件下的热质传递与尺寸收缩模型。采用\textbf{有限体积法}和\textbf{Picard迭代方法}进行数值求解，对药材内部温度、含水率的时空变化及烘干时间进行研究。

\textbf{针对问题一}，构建药材预热阶段的\textbf{径向非稳态热质传递}模型。首先，通过热、质毕渥数判断\textbf{药材内部径向差异}不可忽略；随后，基于圆柱微元的热质通量平衡，建立径向热传导与非线性水分扩散方程，并设置初始和边界条件。采用表面加密有限体积法和\textbf{Crank--Nicolson}格式求解偏微分方程，得到前30 min内温度与含水率的时空分布。1800 s时，中心与表面温度分别为\textbf{33.5753\tempunit}、\textbf{36.7856\tempunit}，含水率分别为\textbf{2.5500\cunit}、\textbf{1.5102\cunit}。结果表明，预热阶段热量传播快于水分迁移，呈现\textbf{整体升温}、\textbf{表层优先失水}的特征。

\textbf{针对问题二}，构建\textbf{变物性热质耦合}干燥模型。在问题一的基础上，进一步考虑密度、比热容和导热系数等参数随含水率的变化，以及水分扩散系数受温度和含水率共同影响，使温度场与水分场形成\textbf{强耦合体系}。采用表面加密有限体积法离散空间，并通过\textbf{中点物性迭代}同步更新温度、含水率及相关物性，得到前3 h内药材温度与含水率的时空分布。3 h时药材温度已基本稳定在\textbf{49.8495}--\textbf{49.9664\tempunit}，而中心与表面含水率仍分别为\textbf{1.7662\cunit}、\textbf{1.0081\cunit}。结果表明，随着干燥进行，温度场趋于均匀，而水分场仍保持明显的内高外低分布，后续干燥逐渐转由内部水分迁移控制。

\textbf{针对问题三}，构建全域含水率约束的\textbf{烘干终点判定}模型。在问题二热质耦合模型的基础上延长计算时间，并以药材各位置含水率均低于0.15\cunit 作为终止条件，逐步跟踪全域最大含水率，并对阈值穿越时刻进行\textbf{插值定位}。计算得到药材整体达到烘干要求的时间为\textbf{57.4721 h}。结果表明，干燥后期水分扩散能力逐渐减弱，\textbf{中心区域成为最后达标位置}，其持续失水决定了整体烘干时长。

\textbf{针对问题四}，构建考虑尺寸收缩的\textbf{动边界热质传递}模型。在问题三的基础上，引入附件给出的半径变化规律，并引入\textbf{物质坐标$\xi=r/R(t)$}，利用其将随时间收缩的物理区域映射到固定区间，从而同步描述尺寸变化与内部水分迁移。沿用有限体积法和中点物性迭代进行数值求解，并以全域含水率低于0.15\cunit 判定烘干终点。计算得到考虑收缩后的整体烘干时间为\textbf{51.0905 h}，较固定尺寸模型\textbf{缩短6.3816 h}。结果表明，尺寸收缩缩短了内部水分迁移距离，使药材后期失水加快，从而提前达到整体烘干要求。

\keywords{药材烘干\quad 径向热质传递\quad 热质耦合\quad 有限体积法\quad 收缩动边界}
\end{abstract}
